\hypertarget{preliminaries}{}
\chapter{Preliminaries: mathematical and statistical foundations}\label{preliminaries}

This chapter assembles, with derivations, every mathematical and
statistical tool that the later chapters use to answer the research
question of Chapter~\ref{the-research-question}. It was built the way the
instructions of a good course demand: by first reading that question and
its complete solution, then listing the technical machinery invoked, and
finally deriving each item from first principles. Section~\ref{prelim-toolkit}
gives that list and maps each topic to the chapter that applies it; the
remaining sections are a self-contained reference. A reader who is fluent
in graduate econometrics may skim and return here whenever a later
chapter cites a result.

Throughout, vectors are columns, a prime or $\top$ denotes transpose,
$\mathbb{E}$, $\operatorname{Var}$ and $\operatorname{Cov}$ are the
expectation, variance and covariance operators, $\phi$ and $\Phi$ are the
standard normal density and cumulative distribution function, and
$I_p$ is the $p\times p$ identity.

\section{The technical toolkit: what is needed and where}\label{prelim-toolkit}

The research question has three parts --- is there benchmark-relative
performance, how uncertain is any ranking of it, and does a past ranking
predict future benchmark-adjusted returns --- and its answer chains
together the following tools. Each row is derived in the section shown and
applied in the chapter shown.

\begin{longtable}[]{@{}p{0.40\textwidth}p{0.15\textwidth}p{0.37\textwidth}@{}}
\toprule\noalign{}
Topic & Derived in & Applied in \\
\midrule\noalign{}
\endhead
\bottomrule\noalign{}
\endlastfoot
Quadratic forms, positive semidefiniteness, spectral theorem & \S\ref{prelim-linalg} & Joint HAC (Ch.~6), shrinkage (Ch.~8), conditioning \\
Matrix inversion lemma and partitioned inverse & \S\ref{prelim-linalg} & Empirical-Bayes posterior (Ch.~8) \\
Multivariate normal; marginal and conditional laws & \S\ref{prelim-mvn} & Empirical Bayes (Ch.~8), GRS (Ch.~7) \\
Quadratic forms of normals; $\chi^2$ and $F$ & \S\ref{prelim-mvn} & Wald and GRS tests (Ch.~7) \\
Ordinary least squares; projection; intercept influence & \S\ref{prelim-ols} & Factor alpha (Ch.~5), HAC (Ch.~6) \\
Asymptotics: LLN, CLT, delta method & \S\ref{prelim-asymp} & All inference \\
Long-run variance; Newey--West HAC & \S\ref{prelim-hac} & Joint alpha covariance (Ch.~6) \\
Wald, GRS, likelihood ratio, boundary mixtures & \S\ref{prelim-tests} & Joint tests (Ch.~7), SFA (Ch.~14) \\
FWER, FDR, Benjamini--Hochberg & \S\ref{prelim-fdr} & Multiplicity control (Ch.~7) \\
Bayes, normal--normal, empirical Bayes, profile MLE & \S\ref{prelim-eb} & Shrinkage (Ch.~8) \\
Bootstrap; block and circular block bootstrap & \S\ref{prelim-boot} & Rank uncertainty (Ch.~9) \\
Spearman correlation; Fisher $z$ & \S\ref{prelim-rank} & Forward evaluation (Ch.~11) \\
Rolling vs.\ expanding; Diebold--Mariano; Giacomini--White & \S\ref{prelim-forecast} & Forward and incremental tests (Ch.~11--12) \\
Composed-error density; conditional mean; boundary LR & \S\ref{prelim-sfa} & Stochastic frontier (Ch.~14) \\
Excess returns and factor models & \S\ref{prelim-ap} & Data and alpha (Ch.~4--5) \\
Numerical linear algebra and optimisation & \S\ref{prelim-num} & Implementation (Ch.~16) \\
\end{longtable}

\section{Linear algebra}\label{prelim-linalg}

\subsection{Quadratic forms and positive semidefiniteness}

For a symmetric matrix $A\in\mathbb{R}^{p\times p}$ and a vector
$x\in\mathbb{R}^{p}$, the scalar $Q(x)=x^{\top}Ax=\sum_{i,j}A_{ij}x_ix_j$
is a \emph{quadratic form}. $A$ is \emph{positive semidefinite} (PSD),
written $A\succeq 0$, if $x^{\top}Ax\ge 0$ for all $x$, and
\emph{positive definite} (PD, $A\succ 0$) if the inequality is strict for
$x\ne 0$. Every covariance matrix is PSD, because for
$\Sigma=\operatorname{Var}(X)$ and any fixed $a$,
\begin{equation}
a^{\top}\Sigma a=\operatorname{Var}(a^{\top}X)\ge 0 .
\end{equation}
This single fact recurs constantly: it is why a HAC covariance must be
PSD to be trusted (Ch.~6), and why the variance of a portfolio combination
of correlated alphas can be smaller than the sum of the individual
variances when the covariance is negative.

\subsection{The spectral theorem, condition number and flooring}

Any symmetric $A$ admits the \emph{eigendecomposition}
\begin{equation}
A=Q\Lambda Q^{\top},\qquad Q^{\top}Q=I_p,\qquad
\Lambda=\operatorname{diag}(\lambda_1,\dots,\lambda_p),
\end{equation}
with orthonormal eigenvectors (columns of $Q$) and real eigenvalues
$\lambda_i$. Then $x^{\top}Ax=\sum_i\lambda_i(q_i^{\top}x)^2$, so
$A\succeq 0$ iff every $\lambda_i\ge 0$. Functions of $A$ act on the
eigenvalues: $A^{-1}=Q\Lambda^{-1}Q^{\top}$ and
$A^{1/2}=Q\Lambda^{1/2}Q^{\top}$ when the $\lambda_i$ are non-negative.
The \emph{condition number} $\kappa(A)=\lambda_{\max}/\lambda_{\min}$
measures how much relative error a linear solve can amplify; the joint HAC
matrix of Chapter~6 has $\kappa\approx 139$, finite and usable.
\emph{Eigenvalue flooring} replaces each $\lambda_i$ by
$\max(\lambda_i,\varepsilon)$ for a small scale-relative $\varepsilon$,
which projects a numerically-indefinite symmetric matrix onto the nearest
PSD matrix in the spectral norm; the implementation floors only negligible
negatives and reports how many, so a materially indefinite estimate raises
rather than being silently repaired.

\subsection{The matrix inversion lemma and the partitioned inverse}

Two identities do the heavy lifting in the empirical-Bayes derivation. The
\emph{Sherman--Morrison--Woodbury} identity is
\begin{equation}\label{eq:woodbury}
(A+UCV)^{-1}=A^{-1}-A^{-1}U\bigl(C^{-1}+VA^{-1}U\bigr)^{-1}VA^{-1},
\end{equation}
valid whenever the inverses exist. It is proved by multiplying the
right-hand side by $A+UCV$ and simplifying to $I$. A consequence used
repeatedly is that, for $V\succ 0$ and $\tau^2>0$,
\begin{equation}\label{eq:eb-woodbury}
\bigl(V^{-1}+\tau^{-2}I\bigr)^{-1}
=\tau^2 I-\tau^4\bigl(V+\tau^2 I\bigr)^{-1},
\end{equation}
which is~\eqref{eq:woodbury} with $A=\tau^{-2}I$, $U=V=I$, $C=V^{-1}$ after
rearrangement; it lets the posterior covariance be written either as a
precision inverse or in the numerically convenient ``$V+\tau^2I$'' form.
For a symmetric block matrix, the \emph{partitioned inverse} uses the
Schur complement $S=A_{22}-A_{21}A_{11}^{-1}A_{12}$:
\begin{equation}
\begin{pmatrix}A_{11}&A_{12}\\A_{21}&A_{22}\end{pmatrix}^{-1}
=\begin{pmatrix}
A_{11}^{-1}+A_{11}^{-1}A_{12}S^{-1}A_{21}A_{11}^{-1} & -A_{11}^{-1}A_{12}S^{-1}\\
-S^{-1}A_{21}A_{11}^{-1} & S^{-1}
\end{pmatrix}.
\end{equation}
This is the algebra behind the conditional normal in \S\ref{prelim-mvn}.

\section{Random vectors and the multivariate normal}\label{prelim-mvn}

\subsection{Expectation, covariance and affine maps}

For a random vector $X$, $\mathbb{E}(X)$ is taken componentwise and
$\operatorname{Cov}(X)=\mathbb{E}[(X-\mathbb{E}X)(X-\mathbb{E}X)^{\top}]$.
The operators are linear under affine maps: if $Y=a+BX$ then
\begin{equation}\label{eq:affine}
\mathbb{E}(Y)=a+B\,\mathbb{E}(X),\qquad
\operatorname{Cov}(Y)=B\operatorname{Cov}(X)B^{\top}.
\end{equation}
Both identities follow directly from linearity of expectation; the second
is the reason a portfolio weight vector applied to correlated returns
produces $w^{\top}\Sigma w$, and the reason the intercept-influence weights
of \S\ref{prelim-ols} turn a residual covariance into an alpha covariance.

\subsection{The multivariate normal density}

$X\sim N(\mu,\Sigma)$ with $\Sigma\succ 0$ has density
\begin{equation}
f(x)=(2\pi)^{-p/2}\lvert\Sigma\rvert^{-1/2}
\exp\!\Bigl(-\tfrac12 (x-\mu)^{\top}\Sigma^{-1}(x-\mu)\Bigr).
\end{equation}
Affine images of normals are normal: by~\eqref{eq:affine},
$a+BX\sim N(a+B\mu,\,B\Sigma B^{\top})$. In particular
$\Sigma^{-1/2}(X-\mu)\sim N(0,I_p)$, the \emph{whitening} transform used
below.

\subsection{Marginal and conditional normals}

Partition $X=(X_1,X_2)$ with mean $(\mu_1,\mu_2)$ and
\begin{equation}
\Sigma=\begin{pmatrix}\Sigma_{11}&\Sigma_{12}\\\Sigma_{21}&\Sigma_{22}\end{pmatrix}.
\end{equation}
Marginals are read off directly: $X_1\sim N(\mu_1,\Sigma_{11})$. The
conditional law is
\begin{equation}\label{eq:condnormal}
X_1\mid X_2=x_2\ \sim\ N\!\Bigl(\mu_1+\Sigma_{12}\Sigma_{22}^{-1}(x_2-\mu_2),\ \Sigma_{11}-\Sigma_{12}\Sigma_{22}^{-1}\Sigma_{21}\Bigr).
\end{equation}
\emph{Derivation.} Write the joint density and collect the terms in $x_1$
in the exponent. Using the partitioned inverse, the precision block acting
on $x_1$ is $(\Sigma_{11}-\Sigma_{12}\Sigma_{22}^{-1}\Sigma_{21})^{-1}$,
and completing the square in $x_1$ shows the conditional exponent is a
quadratic in $x_1$ centred at
$\mu_1+\Sigma_{12}\Sigma_{22}^{-1}(x_2-\mu_2)$ with that precision; the
remaining terms depend only on $x_2$ and form the marginal, confirming the
normalisation. The conditional mean is the linear regression of $X_1$ on
$X_2$ and the conditional covariance is the Schur complement. This is
exactly the shape of the Bayesian posterior in \S\ref{prelim-eb} and of
the truncated-normal update in the frontier model of \S\ref{prelim-sfa}.

\subsection{Quadratic forms of normals; $\chi^2$ and $F$}

If $z\sim N(0,I_p)$ then $z^{\top}z=\sum_i z_i^2\sim\chi^2_p$ by
definition of the chi-square law. Hence for $X\sim N(0,\Sigma)$ with
$\Sigma\succ 0$, whitening gives
\begin{equation}\label{eq:chisq}
X^{\top}\Sigma^{-1}X=(\Sigma^{-1/2}X)^{\top}(\Sigma^{-1/2}X)\sim\chi^2_p .
\end{equation}
This is the engine of the Wald test (\S\ref{prelim-tests}): a standardised
estimate that is asymptotically $N(0,\Sigma)$ under a null produces a
chi-square quadratic form. If $U\sim\chi^2_{d_1}$ and $V\sim\chi^2_{d_2}$
are independent, then $F=(U/d_1)/(V/d_2)\sim F(d_1,d_2)$; the GRS statistic
is exactly such a ratio because it divides a quadratic form in the alphas
by an independent quadratic form in the factor means.

\section{The linear regression model and ordinary least squares}\label{prelim-ols}

\subsection{Model and least-squares objective}

Stack $T$ observations of a portfolio's excess return in $y\in\mathbb{R}^T$
and place a constant and the $K$ factors in the design
$X\in\mathbb{R}^{T\times(K+1)}$, so that with
$\theta=(\alpha,\beta^{\top})^{\top}$,
\begin{equation}
y=X\theta+\varepsilon .
\end{equation}
Ordinary least squares chooses $\theta$ to minimise the residual sum of
squares $S(\theta)=(y-X\theta)^{\top}(y-X\theta)$.

\subsection{Normal equations and the estimator}

Expanding,
$S(\theta)=y^{\top}y-2\theta^{\top}X^{\top}y+\theta^{\top}X^{\top}X\theta$.
Differentiating and setting the gradient to zero,
\begin{equation}
\frac{\partial S}{\partial\theta}=-2X^{\top}y+2X^{\top}X\theta=0
\ \Longrightarrow\ X^{\top}X\hat\theta=X^{\top}y
\ \Longrightarrow\ \hat\theta=(X^{\top}X)^{-1}X^{\top}y,
\end{equation}
the last step assuming $X$ has full column rank. The Hessian
$\partial^2 S/\partial\theta\partial\theta^{\top}=2X^{\top}X\succeq 0$, so
the stationary point is the global minimum.

\subsection{Projection, residuals and the annihilator}

The fitted values are $\hat y=X\hat\theta=Hy$ with the \emph{hat matrix}
$H=X(X^{\top}X)^{-1}X^{\top}$, which is symmetric and idempotent
($H^2=H$) --- the orthogonal projection onto the column space of $X$. The
residuals are $e=(I-H)y=(I-H)\varepsilon$, and $I-H$ is the complementary
projection, so residuals are orthogonal to the regressors, $X^{\top}e=0$.

\subsection{The intercept as a linear combination: influence weights}

Let $c=(1,0,\dots,0)^{\top}$ select the intercept. Then
$\hat\alpha=c^{\top}\hat\theta=c^{\top}(X^{\top}X)^{-1}X^{\top}y$ is a fixed
linear combination of the returns:
\begin{equation}\label{eq:influence}
\hat\alpha=\sum_{t=1}^{T}a_t\,y_t,\qquad
a_t=c^{\top}(X^{\top}X)^{-1}x_t,
\end{equation}
where $x_t^{\top}$ is row $t$ of $X$. Because the model gives
$y_t=\alpha+\beta^{\top}f_t+\varepsilon_t$ and the weights satisfy
$\sum_t a_t=1$, $\sum_t a_t f_t=0$ (the intercept is orthogonal to the
factors by construction), subtracting the truth yields the crucial
identity
\begin{equation}\label{eq:alpha-error}
\hat\alpha-\alpha=\sum_{t=1}^{T}a_t\,\varepsilon_t .
\end{equation}
The estimation error in alpha is a weighted sum of residual shocks, with
weights that depend only on the common factor design. Because every one of
the 25 portfolios shares the same dates and factors, the same weights
$a_t$ apply to all of them, so the same-month residual vector converts
into a joint alpha error --- the observation that produces the full
cross-portfolio covariance in \S\ref{prelim-hac} and Chapter~6.

\subsection{Mean, variance and the sandwich form}

$\hat\theta$ is unbiased, $\mathbb{E}(\hat\theta)=\theta$, and by
\eqref{eq:affine}
\begin{equation}\label{eq:sandwich}
\operatorname{Var}(\hat\theta)=(X^{\top}X)^{-1}X^{\top}\Omega X(X^{\top}X)^{-1},
\qquad \Omega=\operatorname{Var}(\varepsilon\mid X).
\end{equation}
Under the classical assumption $\Omega=\sigma^2 I$ this collapses to
$\sigma^2(X^{\top}X)^{-1}$, and the Gauss--Markov theorem states that OLS is
then the best (minimum-variance) linear unbiased estimator. When $\Omega$
is unknown, heteroskedastic or autocorrelated --- as monthly return
residuals are (Ch.~13) --- the ``bread'' $(X^{\top}X)^{-1}$ is kept but the
``meat'' $X^{\top}\Omega X$ is estimated robustly, giving the HAC estimator
of \S\ref{prelim-hac}.

\section{Asymptotic inference}\label{prelim-asymp}

\subsection{Laws of large numbers and the central limit theorem}

For an ergodic stationary sequence with finite mean, the sample average
converges in probability to the mean (LLN):
$\bar u_T\xrightarrow{p}\mathbb{E}(u)$. If in addition the long-run
variance $\Omega_\infty=\sum_{j=-\infty}^{\infty}\operatorname{Cov}(u_t,u_{t-j})$
is finite, a central limit theorem holds,
\begin{equation}\label{eq:clt}
\sqrt{T}\,(\bar u_T-\mathbb{E}u)\xrightarrow{d}N(0,\Omega_\infty),
\end{equation}
where the variance is the \emph{long-run} variance, not the one-period
variance, precisely because the terms are dependent. This is why forward
evaluation and the Diebold--Mariano test use HAC standard errors on serially
correlated window statistics (\S\ref{prelim-forecast}).

\subsection{Consistency and asymptotic normality of OLS}

Writing $\hat\theta-\theta=(X^{\top}X)^{-1}X^{\top}\varepsilon
=(X^{\top}X/T)^{-1}(X^{\top}\varepsilon/T)$ and letting
$Q=\operatorname*{plim}X^{\top}X/T$, the LLN gives $X^{\top}\varepsilon/T
\xrightarrow{p}0$ so $\hat\theta\xrightarrow{p}\theta$, and the CLT applied
to the score $X^{\top}\varepsilon/\sqrt T$ gives
\begin{equation}
\sqrt T(\hat\theta-\theta)\xrightarrow{d}N\bigl(0,\ Q^{-1}\,S\,Q^{-1}\bigr),
\qquad S=\text{long-run variance of }x_t\varepsilon_t .
\end{equation}
The sandwich \eqref{eq:sandwich} is the finite-sample analogue; $S$ is what
Newey--West estimates.

\subsection{Slutsky and the delta method}

Slutsky's theorem lets consistent estimates of nuisance quantities be
substituted without changing limiting distributions, which is what
justifies plugging an estimated covariance into a Wald statistic. The
\emph{delta method} propagates asymptotic normality through a smooth
function $g$: if
$\sqrt T(\hat\theta-\theta)\xrightarrow{d}N(0,\Sigma)$ then
\begin{equation}\label{eq:delta}
\sqrt T\bigl(g(\hat\theta)-g(\theta)\bigr)\xrightarrow{d}
N\!\bigl(0,\ \nabla g(\theta)^{\top}\Sigma\,\nabla g(\theta)\bigr).
\end{equation}
The Fisher $z$-transform of a correlation (\S\ref{prelim-rank}) is a direct
application.

\section{Dependence-robust covariance estimation}\label{prelim-hac}

\subsection{Autocovariance and the long-run variance}

Let $g_t$ be a mean-zero stationary (vector) score with autocovariances
$\Gamma_j=\mathbb{E}[g_t g_{t-j}^{\top}]$. The variance of the scaled sum
$T^{-1/2}\sum_{t=1}^{T}g_t$ tends to the long-run variance
\begin{equation}\label{eq:lrv}
\Omega_\infty=\sum_{j=-\infty}^{\infty}\Gamma_j
=\Gamma_0+\sum_{j=1}^{\infty}(\Gamma_j+\Gamma_j^{\top}),
\end{equation}
which is $2\pi$ times the spectral density of $g_t$ at frequency zero.
Ignoring the cross terms (using only $\Gamma_0$) is the mistake that
plain OLS standard errors make; when residuals are serially correlated the
$\Gamma_j$ for $j\ge 1$ are non-zero and must be included.

\subsection{The Newey--West / Bartlett estimator}

A truncated, kernel-weighted sample analogue of \eqref{eq:lrv} is the
Newey--West estimator with Bartlett weights,
\begin{equation}\label{eq:nw}
\hat\Omega=\hat\Gamma_0+\sum_{l=1}^{L}w_l\bigl(\hat\Gamma_l+\hat\Gamma_l^{\top}\bigr),
\qquad w_l=1-\frac{l}{L+1},\qquad
\hat\Gamma_l=\frac1T\sum_{t=l+1}^{T}g_t g_{t-l}^{\top}.
\end{equation}
The linearly-declining Bartlett weights are not cosmetic: they make
$\hat\Omega$ positive semidefinite for any data. To see why, the weighted
sum equals $T^{-1}\sum_{|l|\le L}(1-|l|/(L+1))\sum_t g_t g_{t-l}^{\top}$,
which is a Fej\'er-kernel (triangular) average of the periodogram; the
Fej\'er kernel is non-negative, so the resulting matrix is PSD. A sharp
rectangular truncation lacks this property and can yield an indefinite
estimate. The bandwidth $L$ trades bias against variance; a common
data-driven choice is the automatic rule
$L=\lfloor 4(T/100)^{2/9}\rfloor$, used for the forward statistics, while
the monthly factor regressions use a fixed $L=12$.

\subsection{From single-series HAC to the joint alpha covariance}

Combine \eqref{eq:alpha-error} across the $N=25$ portfolios. Let $e_t$ be
the same-month $N$-vector of residuals and let $a_t$ be the common scalar
intercept weight from \eqref{eq:influence}. Define the alpha score vector
$g_t=a_t e_t$. Then the stacked alpha error is
$\hat\alpha-\alpha=\sum_t g_t$, and
\begin{equation}
\operatorname{Var}(\hat\alpha-\alpha)=\sum_{t}\sum_{s}\mathbb{E}[g_t g_s^{\top}]
=\Gamma_0+\sum_{l\ge 1}(\Gamma_l+\Gamma_l^{\top}),
\end{equation}
estimated by \eqref{eq:nw} with the finite-sample correction
$T/(T-P)$, $P=K+1$:
\begin{equation}\label{eq:jointhac}
V_{\mathrm{HAC}}=\frac{T}{T-P}\Bigl\{\hat\Gamma_0+\sum_{l=1}^{L}w_l(\hat\Gamma_l+\hat\Gamma_l^{\top})\Bigr\},
\qquad \hat\Gamma_l=\sum_{t=l+1}^{T}g_t g_{t-l}^{\top}.
\end{equation}
The off-diagonal entries retain contemporaneous and lagged
cross-portfolio dependence; the diagonal reproduces the univariate HAC
variance of each alpha, which is exactly the check performed against
\texttt{statsmodels} in Chapter~6.

\section{Hypothesis testing}\label{prelim-tests}

\subsection{Anatomy of a test}

A test fixes a null hypothesis $H_0$, a statistic whose distribution is
known (at least asymptotically) under $H_0$, and a rejection region. The
\emph{size} is the probability of rejecting a true null; the \emph{power}
is the probability of rejecting a false one. A $p$-value is the
probability, under $H_0$, of a statistic at least as extreme as the one
observed; it is \emph{not} the probability that $H_0$ is true, and a small
$p$-value is evidence against $H_0$ only within the stated model.

\subsection{The Wald test}

To test $H_0:R\theta=r$ for a $q\times(K{+}1)$ matrix $R$, form
\begin{equation}
W=(R\hat\theta-r)^{\top}\bigl[R\hat V R^{\top}\bigr]^{-1}(R\hat\theta-r).
\end{equation}
Under $H_0$, $R\hat\theta-r$ is asymptotically $N(0,R V R^{\top})$, so by
the whitening argument \eqref{eq:chisq}, $W\xrightarrow{d}\chi^2_q$. The
joint test that all 25 alphas are zero is the special case
$R\hat\theta=\hat\alpha$, $r=0$:
\begin{equation}\label{eq:wald}
W=\hat\alpha^{\top}V_{\mathrm{HAC}}^{-1}\hat\alpha\ \xrightarrow{d}\ \chi^2_{25}.
\end{equation}
It is the primary joint test precisely because $V_{\mathrm{HAC}}$ is
dependence-robust: no IID or Gaussian assumption is required for the
limiting chi-square.

\subsection{The $F$ distribution and the GRS test}

Under the stronger classical assumptions --- IID Gaussian residuals with a
constant residual covariance $\Sigma_e$ and factor covariance $\Sigma_f$
with mean $\bar f$ --- the Gibbons--Ross--Shanken statistic has an
\emph{exact} finite-sample $F$ law,
\begin{equation}\label{eq:grs}
\mathrm{GRS}=\frac{T-N-K}{N}\cdot
\frac{\hat\alpha^{\top}\hat\Sigma_e^{-1}\hat\alpha}{1+\bar f^{\top}\hat\Sigma_f^{-1}\bar f}
\ \sim\ F(N,\ T-N-K)\quad\text{under }H_0 .
\end{equation}
It is a ratio of a quadratic form in the alphas to an independent
quadratic form in the factor means, i.e.\ a scaled Hotelling $T^2$, which
is why it is $F$-distributed. Because the diagnostics of Chapter~13 reject
IID normality, GRS is reported only as a conventional benchmark and the
HAC/Wald test is primary.

\subsection{Likelihood ratio, Wilks and boundary mixtures}

For nested models with an unrestricted maximum log-likelihood
$\hat\ell_1$ and a restricted one $\hat\ell_0$, the likelihood-ratio
statistic is $\mathrm{LR}=2(\hat\ell_1-\hat\ell_0)$. Wilks' theorem gives
$\mathrm{LR}\xrightarrow{d}\chi^2_q$ under $H_0$ when the true parameter is
interior to the parameter space. When a tested parameter lies on the
\emph{boundary} --- as in the frontier test of $H_0:\sigma_u=0$ against
$\sigma_u\ge 0$ (\S\ref{prelim-sfa}) --- the standard theory fails, because
the estimator cannot move symmetrically around the null. Under regularity
(Chernoff), the limiting law is the $50{:}50$ mixture
\begin{equation}\label{eq:boundary}
\mathrm{LR}\ \xrightarrow{d}\ \tfrac12\chi^2_0+\tfrac12\chi^2_1,
\qquad p=\tfrac12\Pr(\chi^2_1\ge \mathrm{LR})\ \text{for }\mathrm{LR}>0,
\end{equation}
because with probability one-half the unconstrained estimate of the
non-negative parameter would have been negative and is pinned at zero (the
point mass $\chi^2_0$), and otherwise it behaves as an ordinary
one-degree-of-freedom test.

\section{Multiple testing and false-discovery control}\label{prelim-fdr}

\subsection{Error rates}

Testing $m$ hypotheses inflates the chance of at least one false
rejection. The \emph{family-wise error rate}
$\mathrm{FWER}=\Pr(V\ge 1)$, where $V$ is the number of true nulls
rejected, is controlled conservatively by Bonferroni ($p_i\le\alpha/m$).
The \emph{false-discovery rate} is the expected proportion of false
rejections among all rejections,
\begin{equation}
\mathrm{FDR}=\mathbb{E}\!\left[\frac{V}{\max(R,1)}\right],
\end{equation}
$R$ being the total number of rejections. FDR is the appropriate target
when many tests are screened and a few false positives are tolerable, as
with 25 portfolio alphas.

\subsection{The Benjamini--Hochberg procedure}

Order the $p$-values $p_{(1)}\le\cdots\le p_{(m)}$. Find the largest $k$
with
\begin{equation}
p_{(k)}\le \frac{k}{m}\,q,
\end{equation}
and reject the hypotheses with the $k$ smallest $p$-values. Under
independence or positive dependence (PRDS) this controls the FDR at level
$(m_0/m)q\le q$, where $m_0$ is the number of true nulls. Equivalently one
reports monotone adjusted \emph{q-values}
\begin{equation}\label{eq:bh}
q_{(i)}=\min_{k\ge i}\ \min\!\Bigl(1,\ \frac{m\,p_{(k)}}{k}\Bigr),
\end{equation}
a reverse cumulative minimum that guarantees $q_{(1)}\le\cdots\le q_{(m)}$;
rejecting when $q_i\le q$ reproduces the step-up rule. The study applies
\eqref{eq:bh} in two separate families --- raw-alpha $p$-values and the 25
frontier boundary $p$-values --- because they test different nulls.

\section{Bayesian shrinkage and empirical Bayes}\label{prelim-eb}

\subsection{Bayes' theorem and conjugacy}

Bayes' theorem combines a likelihood and a prior into a posterior,
$p(\theta\mid x)\propto p(x\mid\theta)\,p(\theta)$. A prior is
\emph{conjugate} when the posterior stays in the same family; the normal
likelihood with a normal prior is the canonical case and yields
closed-form shrinkage.

\subsection{The normal--normal posterior by completing the square}

\emph{Scalar case.} Let $x\mid\theta\sim N(\theta,v)$ and
$\theta\sim N(\mu,\tau^2)$. The log-posterior in $\theta$ is, up to
constants,
\begin{equation}
-\tfrac12\Bigl[\frac{(x-\theta)^2}{v}+\frac{(\theta-\mu)^2}{\tau^2}\Bigr]
=-\tfrac12\Bigl[\theta^2\Bigl(\tfrac1v+\tfrac1{\tau^2}\Bigr)
-2\theta\Bigl(\tfrac{x}{v}+\tfrac{\mu}{\tau^2}\Bigr)\Bigr]+\text{const}.
\end{equation}
Completing the square shows $\theta\mid x\sim N(m,s^2)$ with precision
$s^{-2}=v^{-1}+\tau^{-2}$ and
\begin{equation}\label{eq:scalar-post}
m=\frac{v^{-1}x+\tau^{-2}\mu}{v^{-1}+\tau^{-2}}
=\mu+\frac{\tau^2}{\tau^2+v}\,(x-\mu).
\end{equation}
The posterior mean is a precision-weighted average of data and prior; the
\emph{shrinkage weight} $w=\tau^2/(\tau^2+v)$ is near one when the datum is
precise ($v$ small) and near zero when it is noisy.

\emph{Multivariate case.} With $\hat\alpha\mid\alpha\sim N(\alpha,V)$ and
$\alpha\sim N(\mu\mathbf 1,\tau^2 I)$, the same completion of the square in
the vector $\alpha$ gives posterior precision $P=V^{-1}+\tau^{-2}I$ and
\begin{equation}\label{eq:mv-post}
m=\mu\mathbf 1+\tau^2\bigl(V+\tau^2 I\bigr)^{-1}(\hat\alpha-\mu\mathbf 1),
\qquad
C=\tau^2 I-\tau^4\bigl(V+\tau^2 I\bigr)^{-1},
\end{equation}
where the second form of $C$ follows from the Woodbury identity
\eqref{eq:eb-woodbury}. Because $V$ is not diagonal, the update is
genuinely multivariate: one portfolio's posterior depends on uncertainty
it shares with the others.

\subsection{The winner's curse and James--Stein}

Selecting the largest of $N$ noisy estimates biases it upward, because
extreme observed values combine signal with favourable noise. Shrinkage
corrects this: the James--Stein estimator, which shrinks the sample mean
vector toward a common point, has uniformly smaller risk than the raw
estimator in dimension $N\ge 3$. Empirical-Bayes shrinkage of the 25
alphas is the applied form of this dominance and is why the posterior
ranking is more trustworthy than the raw-alpha ranking.

\subsection{Empirical Bayes: estimating $\mu$ and $\tau$}

The prior parameters are not known; empirical Bayes estimates them from
the same cross-section by maximising the \emph{marginal} likelihood.
Integrating $\alpha$ out of the hierarchy gives
\begin{equation}\label{eq:marg}
\hat\alpha\sim N\bigl(\mu\mathbf 1,\ \Sigma_\tau\bigr),
\qquad \Sigma_\tau=V+\tau^2 I .
\end{equation}
For fixed $\tau$, the marginal log-likelihood is quadratic in $\mu$ and its
first-order condition is a generalised-least-squares mean,
\begin{equation}\label{eq:mu-gls}
\frac{\partial\ell}{\partial\mu}=\mathbf 1^{\top}\Sigma_\tau^{-1}(\hat\alpha-\mu\mathbf 1)=0
\ \Longrightarrow\
\hat\mu(\tau)=\frac{\mathbf 1^{\top}\Sigma_\tau^{-1}\hat\alpha}{\mathbf 1^{\top}\Sigma_\tau^{-1}\mathbf 1}.
\end{equation}
Substituting $\hat\mu(\tau)$ back gives a one-dimensional \emph{profile}
log-likelihood in $\tau^2\ge 0$,
\begin{equation}
\ell_p(\tau)=-\tfrac12\log\lvert\Sigma_\tau\rvert
-\tfrac12(\hat\alpha-\hat\mu(\tau)\mathbf 1)^{\top}\Sigma_\tau^{-1}(\hat\alpha-\hat\mu(\tau)\mathbf 1),
\end{equation}
maximised by bounded one-dimensional optimisation (\S\ref{prelim-num}).
The fitted values here are $\hat\mu\approx-26$ and
$\hat\tau\approx 94$ annualised basis points. Reported posterior intervals
add a first-order term for uncertainty in $\hat\mu$ but do not integrate
over $\tau$, so they are model-based summaries, not exact guarantees.

\section{Resampling and the bootstrap}\label{prelim-boot}

\subsection{The plug-in principle}

The bootstrap approximates the sampling distribution of a statistic
$\hat\vartheta=s(\text{data})$ by resampling from an estimate of the
data-generating process and recomputing $s$ many times. For independent
data one resamples observations with replacement; the empirical
distribution of the recomputed statistics approximates its true sampling
distribution. Crucially, the \emph{whole} estimator is recomputed on each
resample, so uncertainty in every stage is captured.

\subsection{Block and circular block bootstraps}

Monthly returns are dependent, so resampling single months would destroy
serial and cross-sectional structure. The \emph{moving block bootstrap}
resamples contiguous blocks of length $b$; concatenating $\lceil T/b\rceil$
blocks preserves dependence within a block. The \emph{circular} variant
wraps the index modulo $T$,
\begin{equation}
\text{block starting at }s:\quad \bigl((s+h-1)\bmod T\bigr)+1,\quad h=0,\dots,b-1,
\end{equation}
so that every observation --- including those near the start and end of the
sample --- is equally likely to appear, removing the edge bias of the
non-circular version. Consistency requires $b\to\infty$ with $b/T\to 0$; a
12-month block is used here. Applying the \emph{same} resampled date index
to all 25 portfolios (a \emph{common-date} bootstrap) preserves
contemporaneous cross-portfolio shocks.

\subsection{Bootstrapping ranks and quantiles}

A rank is a discontinuous function of the scores, so its uncertainty is
not described by a standard error. The bootstrap sidesteps this: within
each replicate $r$ the full pipeline --- regressions, joint HAC covariance,
empirical-Bayes hyperparameters --- is refitted and the portfolios re-ranked
to give $R_i^{(r)}$. The rank interval is the pair of empirical quantiles
\begin{equation}
\bigl[\,Q_{0.025}\{R_i^{(r)}\},\ Q_{0.975}\{R_i^{(r)}\}\,\bigr],
\end{equation}
and the top-quintile probability is the fraction of replicates with
$R_i^{(r)}\le 5$. These are resampling distributions, not Gaussian
intervals, and can be asymmetric and bounded at ranks $1$ and $25$.

\section{Rank statistics}\label{prelim-rank}

\subsection{Spearman rank correlation}

Spearman's $\rho$ is the Pearson correlation computed on the ranks of two
variables; it measures monotone association and is invariant to any
increasing transformation of either variable. In the forward evaluation it
compares the cross-section of training scores with the cross-section of
future alphas, so that only the \emph{ordering}, not the scale, is tested.

\subsection{The Fisher $z$-transformation}

A sample correlation near $\pm 1$ has a strongly skewed, bounded sampling
distribution, so correlations are averaged on the \emph{Fisher} scale
\begin{equation}
z=\operatorname{atanh}(\rho)=\tfrac12\log\frac{1+\rho}{1-\rho}.
\end{equation}
This is variance-stabilising. By the delta method \eqref{eq:delta} with
$g(\rho)=\operatorname{atanh}(\rho)$, $g'(\rho)=1/(1-\rho^2)$, and the
approximate correlation variance $\operatorname{Var}(\hat\rho)\approx
(1-\rho^2)^2/n$,
\begin{equation}
\operatorname{Var}(z)\approx g'(\rho)^2\operatorname{Var}(\hat\rho)\approx\frac{1}{n},
\end{equation}
free of $\rho$ to first order. Window correlations are therefore averaged
as $z$, given HAC standard errors across windows, and mapped back with
$\tanh$.

\section{Out-of-sample forecast evaluation}\label{prelim-forecast}

\subsection{Rolling vs.\ expanding estimation and look-ahead bias}

A \emph{rolling} scheme estimates on a fixed-length trailing window
$[t-w+1,t]$ and predicts the next $h$ months; an \emph{expanding} scheme
uses all data $[1,t]$. \emph{Look-ahead bias} is any leakage of information
dated after $t$ into the time-$t$ prediction --- a future return entering a
score, rank, prior parameter or factor loading. The study freezes every
such quantity at the training cutoff and evaluates only strictly later
months, and asserts the non-overlap of evaluation windows so that annual
statistics are not mechanically correlated through shared future data.

\subsection{The Diebold--Mariano test}

To compare two forecasts with losses $L(e_{1t})$ and $L(e_{2t})$, form the
loss differential $d_t=L(e_{1t})-L(e_{2t})$ and test $H_0:\mathbb{E}(d_t)=0$
(equal predictive accuracy) with
\begin{equation}\label{eq:dm}
\mathrm{DM}=\frac{\bar d}{\sqrt{\widehat{\operatorname{Var}}_{\mathrm{HAC}}(\bar d)}}
\ \xrightarrow{d}\ N(0,1),
\end{equation}
where the HAC variance \eqref{eq:nw} is needed because $d_t$ is serially
correlated. The incremental test of Chapter~12 is exactly this: $d_t$ is
the per-window difference between the rolling ranking's predictive score
(rank correlation, or quintile spread) and a benchmark ordering's, and the
reported HAC $t$-statistic on $\bar d$ is a Diebold--Mariano statistic.

\subsection{Giacomini--White conditional predictive ability}

The classical Diebold--Mariano theory treats forecasts as given.
Giacomini and White extend equal-accuracy testing to
forecasts produced by models that are \emph{re-estimated} on a rolling
window of fixed length, which is precisely the design here; their result
justifies treating the rolling-versus-benchmark differential as a valid
test of conditional predictive ability rather than merely of two fixed
forecasts.

\section{Composed-error (stochastic frontier) models}\label{prelim-sfa}

\subsection{The composed error and its density}

The half-normal stochastic-frontier model writes the residual as
$\varepsilon=v-u$ with $v\sim N(0,\sigma_v^2)$ and a one-sided
$u=\lvert w\rvert$, $w\sim N(0,\sigma_u^2)$, independent. Its density is the
convolution
\begin{equation}
f(\varepsilon)=\int_0^{\infty} f_v(\varepsilon+u)\,f_u(u)\,du
=\frac{2}{\sigma}\,\phi\!\Bigl(\frac{\varepsilon}{\sigma}\Bigr)\,
\Phi\!\Bigl(-\frac{\lambda\varepsilon}{\sigma}\Bigr),
\qquad \sigma^2=\sigma_v^2+\sigma_u^2,\ \lambda=\frac{\sigma_u}{\sigma_v}.
\end{equation}
\emph{Derivation.} The integrand is
$\propto\exp\{-\tfrac12[(\varepsilon+u)^2/\sigma_v^2+u^2/\sigma_u^2]\}$;
completing the square in $u$ shows $u$ enters as a normal kernel with mean
$\mu_*=-\varepsilon\sigma_u^2/\sigma^2$ and variance
$\sigma_*^2=\sigma_v^2\sigma_u^2/\sigma^2$, and the leftover $\varepsilon$
terms assemble into $\phi(\varepsilon/\sigma)$; integrating the normal
kernel over $u\ge 0$ contributes the factor
$\Phi(\mu_*/\sigma_*)=\Phi(-\lambda\varepsilon/\sigma)$, giving the stated
density (the factor $2$ is the half-normal normalisation).

\subsection{The conditional mean $\mathbb{E}[u\mid\varepsilon]$}

From the completion above, $u\mid\varepsilon$ is normal with mean $\mu_*$
and variance $\sigma_*^2$ \emph{truncated} to $[0,\infty)$. The mean of a
normal $N(\mu_*,\sigma_*^2)$ truncated below at zero is the Jondrow et
al.\ formula
\begin{equation}\label{eq:jondrow}
\mathbb{E}[u\mid\varepsilon]=\mu_*+\sigma_*\,
\frac{\phi(\mu_*/\sigma_*)}{\Phi(\mu_*/\sigma_*)},
\end{equation}
the standard truncated-normal mean (a location plus the inverse Mills ratio
scaled by $\sigma_*$). This is the portfolio-specific estimate of the
one-sided shortfall used, cautiously, as a residual-asymmetry diagnostic.

\subsection{The boundary likelihood-ratio test}

Testing whether a distinct one-sided component exists is $H_0:\sigma_u=0$,
a value on the boundary of the admissible set $\sigma_u\ge 0$. As in
\eqref{eq:boundary}, the likelihood-ratio statistic against the Gaussian
model has the $\tfrac12\chi^2_0+\tfrac12\chi^2_1$ limiting law, so the
$p$-value is $\tfrac12\Pr(\chi^2_1\ge\mathrm{LR})$. Only one of the 25
portfolios survives this test after false-discovery control, which is why
the frontier apparatus remains a narrow diagnostic rather than a ranking.

\section{Asset-pricing setup}\label{prelim-ap}

\subsection{What ``Fama--French'' means}\label{prelim-ff}

The names attached to the model and the data are those of the financial
economists Eugene F.\ Fama (a Nobel laureate in economics, 2013) and
Kenneth R.\ French, whose joint work defines the benchmark used throughout
these notes. Understanding the term requires one step of background.

\emph{The starting point: the CAPM.} The classical Capital Asset Pricing
Model (Sharpe, Lintner) holds that a single systematic factor --- the
excess return of the aggregate market --- prices all assets: an asset's
expected excess return should be proportional to its market ``beta,''
\begin{equation}
\mathbb{E}[y_{i}]=\beta_{i}\,\mathbb{E}[\text{Mkt}-\text{RF}],
\end{equation}
with intercept zero. Empirically this fails in a systematic way. Two
groups of stocks earn persistently higher average returns than their
market beta explains: \emph{small} firms (low market capitalisation) and
\emph{value} firms (high book-to-market ratio). These are the ``size'' and
``value'' anomalies.

\emph{The three-factor model (FF3).} Fama and French (1993) proposed
augmenting the market factor with two \emph{factor-mimicking portfolios}
built precisely to capture those anomalies: SMB (``small minus big,'' long
small firms and short large ones) and HML (``high minus low,'' long value
and short growth). ``FF3'' abbreviates the resulting Fama--French
three-factor model. It has been the standard empirical risk benchmark in
finance for three decades, which is exactly why it is the benchmark here:
alpha measured against FF3 is the conventional definition of return
``unexplained'' by known systematic exposures.

\emph{The five-factor extension (FF5).} Fama and French (2015) added two
further factors --- RMW (``robust minus weak'' profitability) and CMA
(``conservative minus aggressive'' investment) --- to capture patterns FF3
still missed. ``FF5'' is that five-factor model. The study uses FF3 as its
canonical benchmark and FF5 as a prespecified sensitivity check, precisely
because alpha is benchmark-dependent (\S\ref{rq-difficulties}).

\emph{The test assets and the data source.} The 25 portfolios analysed
here are themselves a Fama--French construction: all NYSE/AMEX/NASDAQ
stocks are sorted, at each formation date, into five size groups and five
book-to-market groups, and the $5\times5=25$ intersections are held as
value-weighted portfolios. Both the factor series and the 25 portfolios
are published and maintained in the \emph{Kenneth R.\ French Data Library},
the canonical public source from which the local inputs are drawn
(Chapter~4). A subtle but important consequence follows: because SMB and
HML are built \emph{from} sorts on size and book-to-market, and the 25 test
assets are built from the \emph{same} sorts, regressing these portfolios on
FF3 measures the model against its own raw material --- another reason a
non-zero alpha is a pricing residual, not evidence of skill.

\subsection{Excess returns and factor models}

The one-period excess return is $y_{it}=r_{it}-r_{ft}$, the return above
the risk-free rate. A linear factor model writes it as
$y_{it}=\alpha_i+\beta_i^{\top}f_t+\varepsilon_{it}$, where $f_t$ collects
traded factor returns. For the three-factor model
$f_t=(\text{Mkt}-\text{RF},\,\text{SMB},\,\text{HML})_t^{\top}$: the market
excess return, SMB (small minus big) and HML (high minus low
book-to-market); the five-factor model appends RMW (profitability) and
CMA (investment). Every component is itself an excess or long--short return,
so the regression places dependent variable and regressors on a common
basis.

\subsection{Alpha as a model-relative residual}

The intercept $\alpha_i$ is the average excess return the chosen factors
leave unexplained. It is model-relative: a non-zero alpha signals mispricing
\emph{relative to the specified factors}, and can arise from an omitted
priced risk, a friction, a nonlinearity, or genuine abnormal performance.
On passively sorted test assets it is a pricing residual, not evidence of
skill --- the interpretive discipline that the whole course enforces.

\section{Numerical methods}\label{prelim-num}

\subsection{Solve, do not invert}

Equations are written with matrix inverses for clarity, but a stable
implementation solves linear systems instead. For a PD matrix the Cholesky
factorisation $A=LL^{\top}$ ($L$ lower-triangular) reduces $Ax=b$ to two
triangular solves, at lower cost and better accuracy than forming
$A^{-1}$. The log-determinant needed for the marginal likelihood
\eqref{eq:marg} is $\log\lvert A\rvert=2\sum_i\log L_{ii}$.

\subsection{Conditioning and eigenvalue flooring}

The relative error in a solve is bounded by roughly $\kappa(A)$ times the
input error, so an ill-conditioned covariance loses precision. The pipeline
symmetrises every estimated covariance, inspects its eigenvalues, floors
only negligible negative eigenvalues by a scale-relative amount
(\S\ref{prelim-linalg}), and raises if the matrix is materially indefinite,
so numerical repair is auditable rather than silent.

\subsection{Positivity and bounded optimisation}

Scale parameters must stay positive, so they are estimated on the log
scale, e.g.\ $\sigma=e^{\eta}$ with $\eta$ unconstrained, which turns a
constrained problem into an unconstrained one and keeps the optimiser
inside the feasible region. The one-dimensional profile in $\tau^2$
\eqref{eq:mu-gls} is maximised by a bounded scalar routine (Brent's method)
on $[0,\text{upper}]$, and the frontier likelihood uses a robust
Newton/quasi-Newton step with the analytic score of Chapter~14.

